\documentclass{report}
\usepackage[utf8]{inputenc}
\usepackage[T1]{fontenc}
\usepackage{geometry}
\usepackage{graphicx}
\usepackage{float}
\usepackage{hyperref}
\usepackage[english,indonesian]{babel}

\title{Optimasi Pemrosesan Template LaTeX Menggunakan Pendekatan Berbasis AI}
\author{Academic Support Agent}
\date{\today}

\begin{document}
\maketitle

\begin{abstract}

\end{abstract}

\chapter{PENDAHULUAN}


\chapter{TINJAUAN PUSTAKA}


% Original content (commented out):
% %NO BORRAR PARA COMPILAR PREGUNTA%%%%%%
\ifdefined\niveldos\else
\documentclass[12pt,letterpaper]{report}
\usepackage{DefinicionClases}
\begin{document}
\fi
%%%%%%%%%%%%%%%%%%%%%%%%%%%%%%%%%%%%%%%%%
%%%%%%%%%%%%%%%%%%%%%%%%%%%%%%%%%%%%%%%%%%%%%%%%%%%%%%%%%%
\claseaclase{A5}

\begin{locro}
\item  Resolver problemas en coordenadas cilíndricas
\end{locro}

\begin{contenidos}
\item Aplicaciones de clase A4
\end{contenidos}


\begin{actividad}

Clase expositivas 

Se desarrollan los problemas propuestos.
\end{actividad}

%%%%%%%%%%%%%%%%%%%%%%%%%%%%%%%%%%%%%%%%%%%%%%%%%%%%%%%%%%

\begin{ejemploclase}{}

 Una partícula se mueve en el plano $XY$ con \underline{rapidez constante} $V_0$ a lo largo de una curva dada por la ecuación $r=k(1+\cos \theta)$, con $k$ constante positiva. Esta curva es conocida como cardioide.

\begin{enumerate}
\item Calcule la velocidad angular de la part\'icula en función del ángulo $\theta$. 
\item Pruebe que la aceleración radial de la partícula es constante y calcule su valor.
\end{enumerate}

\end{ejemploclase} 


\begin{solucion}
  La ecuación del cardioide está dada por $r=k(1+\cos \theta)$ por
lo tanto:

\[ \dot{r}= -k \dot\theta \sin\theta \]
\[ \ddot{r}= -k (\dot{\theta}^2 \cos\theta +\ddot{\theta}
\sin\theta)\]



 En coordenadas polares la velocidad de la partícula está dada por
la expresión:

\[ \vec{v}= \dot{r} \hat{r}+r\dot{\theta} \hat{\theta}\]

 Como la rapidez de la partícula es constante entonces

\[ V_0^2 = \dot{r}^2+r^2\dot{\theta}^2=2 k^2 \dot{\theta}^2
(1+\cos\theta)\]

\[\Longrightarrow
\boxed{\dot{\theta}=\frac{V_0}{\sqrt{2k^2(1+\cos\theta)}}=\frac{V_0}{\sqrt{2kr}}}\]

 Para calcular la aceleración radial $a_r=\ddot{r}-r\dot{\theta}^2$
necesitamos calcular $\ddot{\theta}$.

\[\dot{\theta}=\frac{V_0}{\sqrt{2kr}}\Longrightarrow \ddot{\theta}=-\frac{V_0 \dot{r}}{2r\sqrt{2kr}}=\frac{k\sin\theta}{2r}\dot{\theta}^2=\frac{V_0^2 \sin\theta}{4r^2}\]

 Luego

\[a_r=\ddot{r}-r\dot{\theta}^2= -k\cos\theta
\dot{\theta}^2-k\sin\theta\ddot{\theta}-r\dot{\theta}^2\]
\[\Longrightarrow a_r = -k\dot{\theta}^2\left[
\cos\theta+\frac{\sin^2\theta}{2(1+\cos\theta)}+(1+\cos\theta)\right]\]
\[a_r=-k\dot{\theta}^2\left[\frac{2\cos\theta+2\cos^2\theta+\sin^2\theta+2+4\cos\theta+2\cos^2\theta}{2(1+\cos\theta)}
\right]\]
\[a_r=\frac{-k\dot{\theta}^2}{2(1+\cos\theta)}(4\cos^2+\sin^2\theta+6\cos\theta+2))= \frac{-3k\dot{\theta}^2}{2(1+\cos\theta)}(1+\cos\theta)^2\]

\[ \boxed{a_r= -\frac{3}{4}\frac{V_0^2}{k}} \]

\end{solucion} 

%%%%%%%%%%%%%%%%%%%%%%%%%%%%%%%%%%%%%%%%%%%%%%%%%%%%%%%%%%


\begin{ejemploclase}{FMF122_07_clase_A5_01}

 La trayectoria de vuelo del avión $B$ es una línea horizontal que pasa exactamente por la vertical de la estación de radar $A$. Si el avión viaja hacia la izquierda con velocidad constante $v_0$. Determine:
 
\begin{enumerate}
\item La velocidad angular del radar ($\dot{\theta}$) en función de $v_0$, $h$ y $\theta$.
\item  La aceleración angular del radar ($\ddot{\theta}$) en función de $v_0$, $h$ y $\theta$.
\end{enumerate}

\vspace{30mm}

\end{ejemploclase} 


%%%%%%%%%%%%%%%%%%%%%%%%%%%%%%%%%%%%%%%%%%%%%%%%%%%%%%%%%%


\begin{ejemploclase}[0.6]{FMF122_07_clase_A5_01}

 La trayectoria de vuelo del avión $B$ es una línea horizontal que pasa exactamente por la vertical de la estación de radar $A$. Si el avión viaja hacia la izquierda con velocidad constante $v_0$. Determine:
 
\begin{enumerate}
\item La velocidad angular del radar ($\dot{\theta}$) en función de $v_0$, $h$ y $\theta$.
\item  La aceleración angular del radar ($\ddot{\theta}$) en función de $v_0$, $h$ y $\theta$.
\end{enumerate}

\vspace{60mm}

\end{ejemploclase} 


%%%%%%%%%%%%%%%%%%%%%%%%%%%%%%%%%%%%%%%%%%%%%%%%%%%%%%%%%%

%%%%%%%%%%%%%%%%%%%%%%%%%%%%%%%%%%%%%%%%%%%%%%%%%%%%%%%%%%


\begin{tareacasa}{FMF122_07_clase_A5_01}

 La trayectoria de vuelo del avión $B$ es una línea horizontal que pasa exactamente por la vertical de la estación de radar $A$. Si el avión viaja hacia la izquierda con velocidad constante $v_0$. Determine:
 
\begin{enumerate}
\item La velocidad angular del radar ($\dot{\theta}$) en función de $v_0$, $h$ y $\theta$.
\item  La aceleración angular del radar ($\ddot{\theta}$) en función de $v_0$, $h$ y $\theta$.
\end{enumerate}

\vspace{30mm}

\end{tareacasa} 


%%%%%%%%%%%%%%%%%%%%%%%%%%%%%%%%%%%%%%%%%%%%%%%%%%%%%%%%%%


\begin{tareacasa}[0.6]{FMF122_07_clase_A5_01}

 La trayectoria de vuelo del avión $B$ es una línea horizontal que pasa exactamente por la vertical de la estación de radar $A$. Si el avión viaja hacia la izquierda con velocidad constante $v_0$. Determine:
 
\begin{enumerate}
\item La velocidad angular del radar ($\dot{\theta}$) en función de $v_0$, $h$ y $\theta$.
\item  La aceleración angular del radar ($\ddot{\theta}$) en función de $v_0$, $h$ y $\theta$.
\end{enumerate}

\vspace{30mm}

\end{tareacasa} 




%NO BORRAR PARA COMPILAR PREGUNTA%%%%%%%%%
\ifdefined\niveldos\else
\end{document} 
\fi
%%%%%%%%%%%%%%%%%%%%%%%%%%%%%%%%%%%%%%%%%%


\end{document}